%% ============================================================
%%  Chapitre 7 — Analytics et Preuve d'Edge
%% ============================================================
\chapter{Analytics et Résultats}
\label{ch:analytics}

\section{Introduction}
\label{sec:intro_analytics}

Après avoir présenté les moteurs de signaux (Chapitres~\ref{ch:signaux} et \ref{ch:wfo}),
l'architecture de la plateforme (Chapitre~\ref{ch:architecture}) et le tableau de bord
décisionnel (Chapitre~\ref{ch:dashboard}), il est impératif de prouver rigoureusement
la valeur prédictive des signaux générés.

Ce chapitre se concentre sur le module \textbf{Analytics}, conçu pour fournir aux
utilisateurs (traders, analystes, jury) une preuve statistique claire et indiscutable
de la présence d'un \textit{edge} (avantage prédictif) sur le marché.

\section{Analyse par Titre : Return Attendu et Hit Rate}

La vue analytique par action offre une évaluation directe de la qualité prédictive
d'un signal. Pour chaque niveau de signal généré (Achat Fort, Achat, Neutre, Vente, Vente Forte),
le système calcule :

\subsection{L'Expected Return (Rendement Attendu)}
Le rendement moyen attendu est calculé pour différents horizons prédictifs
(\textit{forward returns} à 1j, 5j et 21j). Pour que la mesure soit statistiquement
valide, un intervalle de confiance (CI) à 95\,\% est fourni par \textit{bootstrap}.

Afin de s'assurer que le rendement attendu n'est pas le fruit du hasard, un test
de Monte Carlo non paramétrique (Luck Test) est appliqué. Les rendements sont centrés~:
\begin{equation}
\tilde{r}_t = r_t - \mu_r
\end{equation}
Puis, $N=2000$ trajectoires factices sont générées par tirage aléatoire avec remise.
La p-value empirique correspond à la fraction des trajectoires dont le rendement
moyen simulé dépasse le rendement moyen observé.

\subsection{Le Hit Rate}
Le \textit{hit rate} représente la probabilité historique que le signal donne lieu à
un rendement strictement positif à l'horizon considéré. Un \textit{hit rate}
significativement supérieur à 50 \% (ou inférieur à 50 \% pour un signal de vente)
démontre une asymétrie exploitable. Comme pour l'expected return, il s'accompagne
d'un intervalle de confiance.

\section{Preuve d'Edge : La Matrice de Capacité Prédictive}

La preuve absolue de l'edge d'un système de signaux réside dans la \textbf{monotonie}
des résultats. La matrice de capacité prédictive (Predictive Ability Matrix) croise les
labels de signaux avec les horizons forward.

Un système robuste présente une structure monotone :
\begin{center}
$E[R \,|\, \text{Achat Fort}] > E[R \,|\, \text{Achat}] > E[R \,|\, \text{Neutre}] > E[R \,|\, \text{Vente}] > E[R \,|\, \text{Vente Forte}]$
\end{center}

Cette monotonie, mesurée via le \textit{Score monotonicity},
prouve que la force du signal (son score d'ensemble pondéré) est proportionnelle
au mouvement futur des prix :
\begin{equation}
\text{Monotonicity} = E[R \,|\, \text{Achat Fort}] - E[R \,|\, \text{Vente Forte}]
\end{equation}

Le tableau~\ref{tab:predictive_matrix} illustre un exemple de matrice
de capacité prédictive typique extraite de la plateforme sur une valeur du MASI
(validation sur sous-périodes récentes). L'Expected Return s'apprécie à mesure de la force
du signal.

\begin{table}[H]
\centering
\caption{Exemple de Matrice de Capacité Prédictive (Rendements forward annualisés)}
\label{tab:predictive_matrix}
\begin{tabularx}{\linewidth}{@{}lYYYY@{}}
\toprule
\textbf{Bucket (Signal)} & \textbf{N (Obs)} & \textbf{1 jour} & \textbf{5 jours} & \textbf{21 jours} \\
\midrule
Achat Fort ($> 50$)         & 142 & +14.2\,\% & +11.8\,\% & +8.5\,\% \\
Achat ($15$ à $50$)         & 381 &  +6.1\,\% &  +4.2\,\% & +3.9\,\% \\
Neutre ($-15$ à $15$)       & 914 &  +0.5\,\% &  -0.1\,\% & -0.4\,\% \\
Vente ($-50$ à $-15$)       & 342 &  -4.8\,\% &  -3.5\,\% & -2.1\,\% \\
Vente Forte ($< -50$)       & 128 & -11.5\,\% &  -9.2\,\% & -6.8\,\% \\
\midrule
\textbf{Monotonicity Score} & & \textbf{+25.7\,\%} & \textbf{+21.0\,\%} & \textbf{+15.3\,\%} \\
\bottomrule
\end{tabularx}
\end{table}

\section{Information Coefficient (IC) et Preuves Avancées}

\subsection{Les Courbes IC (Information Coefficient)}
Le coefficient d'information de Spearman (Rank IC) mesure la corrélation entre les
scores des signaux et les rendements futurs. Les courbes IC à 1j, 2j, 3j, 5j et 10j
sont générées avec des intervalles de confiance bootstrap. Un IC stable et positif
confirme le pouvoir de prédiction du signal.

\subsection{Rigueur Scientifique : DSR et Haircut}
Pour répondre définitivement au problème des tests multiples, la plateforme inclut~:
\begin{itemize}
    \item \textbf{Deflated Sharpe Ratio (DSR)}~: Ajuste la significativité du Sharpe
    en fonction du nombre de variantes testées, tenant compte de la variance de
    l'échantillon des Sharpe observés (Bailey \& López de Prado, 2014).
    \item \textbf{Harvey-Liu Sharpe Haircut}~: Pénalité appliquée au Sharpe évalué sur
    les sous-périodes pour corriger le biais de sélection (Harvey, Liu \& Zhu, 2016). Si $\widehat{\text{SR}}$
    est le Sharpe observé, le "haircut" est calculé par~:
    \begin{equation}
    \text{Haircut} = \max\left(0, \; 1 - \frac{E[\text{SR}_{\max}(N)]}{\widehat{\text{SR}}}\right)
    \end{equation}
    où $E[\text{SR}_{\max}(N)]$ est l'espérance du maximum de $N$ tirages normaux (nombre de variantes).
    \item \textbf{Probabilistic Sharpe Ratio (PSR)}~: Ajuste l'estimation du Sharpe
    en fonction de l'asymétrie (\textit{skewness}) et de la voussure (\textit{kurtosis}) des rendements,
    prouvant que la performance ne repose pas sur un artefact de la loi normale.
\end{itemize}

\subsection{Contrôle du False Discovery Rate (FDR)}
Sur la batterie globale d'indicateurs (les familles testées sur plusieurs horizons),
le taux de fausses découvertes est strictement contrôlé par la procédure de
\textbf{Benjamini-Hochberg (1995)}.

Les p-values $p_i$ de chaque signal sont triées en ordre croissant $p_{(1)} \le p_{(2)} \le \dots \le p_{(m)}$.
Le système rejette l'hypothèse nulle pour toutes les hypothèses $i \le k$,
où $k$ est le plus grand indice vérifiant la limite de tolérance $q = 0.10$~:
\begin{equation}
p_{(k)} \le \frac{k}{m} \cdot q
\end{equation}

L'intégration native de ce filtre garantit que le pipeline ne remonte au desk
de trading que des signaux statistiquement robustes, minimisant le risque
d'overfitting décrit dans la problématique.

\section*{Conclusion du chapitre}

Le module Analytics constitue la clé de voûte scientifique de la plateforme. En
fournissant des intervalles de confiance rigoureux, en prouvant la monotonie de l'edge,
et en appliquant les ajustements avancés (DSR/PSR), la plateforme garantit à l'utilisateur
que les signaux présentés sur le Dashboard possèdent une validité robuste et répétable,
répondant ainsi avec précision à la problématique initiale.